\documentclass[12pt]{article}
\usepackage[margin=0.58in,top=0.62in,bottom=0.48in]{geometry}
\usepackage{lmodern}
\usepackage{microtype}
\usepackage{fancyhdr}
\usepackage{titlesec}
\usepackage{enumitem}
\usepackage{lastpage}
\usepackage[hidelinks]{hyperref}

\setlength{\parindent}{0pt}
\setlength{\parskip}{0.18em}
\setlength{\headheight}{26pt}
\raggedbottom
\titleformat{\section}{\normalsize\bfseries\scshape}{}{0pt}{}
\titleformat{\subsection}{\normalsize\bfseries}{}{0pt}{}
\titlespacing*{\section}{0pt}{0.35em}{0.12em}
\titlespacing*{\subsection}{0pt}{0.25em}{0.08em}
\pagestyle{fancy}
\fancyhf{}
\fancyhead[C]{\footnotesize Lit Example Reader \quad Assignment ID: U1L14LE \quad Page \thepage\ of \pageref{LastPage}}
\renewcommand{\headrulewidth}{0.5pt}

\newenvironment{playpassage}{\begin{quote}\footnotesize\setlength{\parskip}{0.08em}\setlength{\parindent}{0pt}}{\end{quote}}
\newcommand{\speaker}[1]{\textbf{#1.}}

\begin{document}

\begin{center}
{\LARGE\bfseries Week 14 Lit Example Reader}\par\vspace{0.02em}
{\large\scshape Shakespeare: Discovery, Proof, and Judgment}\par\vspace{0.06em}
\rule{0.78\textwidth}{0.5pt}
\end{center}

\textbf{Purpose:} Shakespeare passages for Week 14: ask what had to be discovered, what was wrongly inferred, and what judgment failed before the discovery arrived.

\textbf{What to watch for:} Characters often have words, signs, predictions, or warnings before they have the full facts. Notice where they need observation, testimony, clarified terms, or better judgment.

\section*{Before the Passages: The Week 14 Logic Tool}

\begin{itemize}[leftmargin=1.3em,itemsep=0.05em,topsep=0.08em,parsep=0.03em]
\item \textbf{Discovery:} learning a fact not yet known.
\item \textbf{Proof:} showing why a conclusion follows from reasons.
\item \textbf{Judgment:} choosing the right tool and sizing the conclusion to the evidence.
\item \textbf{Tool map:} question; needed discovery; clarified terms; inference test; judgment.
\end{itemize}

\section*{Passage 1: Macbeth, Act 5 --- Discovering the Meaning Too Late}

\textbf{Context:} Macbeth trusted the apparitions' words: Birnam Wood would come to Dunsinane, and no one born of woman would harm him. In Act 5, he begins to discover what the words really allowed.

\begin{playpassage}
\speaker{Messenger} As I did stand my watch upon the hill,\\
I looked toward Birnam, and anon, methought,\\
The wood began to move.\\
\speaker{Macbeth} If thou speak'st false,\\
Upon the next tree shalt thou hang alive,\\
Till famine cling thee: if thy speech be sooth,\\
I care not if thou dost for me as much.\\
...\\
\speaker{Macduff} Despair thy charm;\\
And let the angel whom thou still hast served\\
Tell thee, Macduff was from his mother's womb\\
Untimely ripp'd.\\
\speaker{Macbeth} Accursed be that tongue that tells me so,\\
For it hath cow'd my better part of man!
\end{playpassage}

\subsection*{Reader Notes}

This is your assisted example. Macbeth had testimony from the apparitions, but he judged the phrases too quickly. He treated strange words as if their meaning was already settled. The moving wood is not a forest walking by magic; it is an army using branches. ``Born of woman'' did not cover Macduff in the ordinary way Macbeth assumed. The late discovery exposes a failure of judgment: Macbeth inferred safety before he had clarified the terms and tested the source.

\textbf{Reading task:} Mark two facts Macbeth discovers too late. In the margin, write which tool was needed earlier: clarified terms, source testing, observation, or judgment.

\newpage

\section*{Passage 2: Julius Caesar, Act 2 --- Signs, Warnings, and Judgment}

\textbf{Context:} Calpurnia urges Caesar to stay home because of terrifying signs and dreams. Decius then reinterprets the dream in a flattering way. Caesar must judge what kind of evidence he has.

\begin{playpassage}
\speaker{Calpurnia} Caesar, I never stood on ceremonies,\\
Yet now they fright me. There is one within,\\
Besides the things that we have heard and seen,\\
Recounts most horrid sights seen by the watch.\\
...\\
\speaker{Caesar} The things that threaten'd me\\
Ne'er look'd but on my back; when they shall see\\
The face of Caesar, they are vanished.\\
...\\
\speaker{Decius} This dream is all amiss interpreted;\\
It was a vision fair and fortunate.\\
Your statue spouting blood in many pipes,\\
In which so many smiling Romans bathed,\\
Signifies that from you great Rome shall suck\\
Reviving blood.
\end{playpassage}

\subsection*{Reader Notes}

This passage shifts more work to you. Several kinds of material appear: Calpurnia's fear, reported sights, Caesar's confidence, and Decius's flattering interpretation. Do not decide simply by asking which speech sounds stronger. Ask what is actually known, what is merely interpreted, and what judgment Caesar needs before acting.

\textbf{Reading task:} Underline one warning, circle one interpretation, and put a star beside one line where pride or pressure may distort Caesar's judgment. Save the final tool map for the worksheet.

\section*{How to Use These Passages on the Worksheet}

Do not retell the plot. Use Week 14's tool map: question, needed discovery, clarified term or sign, inference test, and judgment.

\end{document}
